%%% FIELD def-effect-image
Let \(e_u\in\mathbb R^V\) be the standard basis vector at label \(u\).  Define the joint intervention-column image by
\[
S_{\bullet u}(W,\mathcal G,\ell)=
\begin{cases}
\{e_u\},&g_{\mathcal G,\ell}(u)=\bot,\\
\displaystyle
\left\{\frac{W_{\bullet c}}{W_{uc}}:(u,c)\in
P_{g_{\mathcal G,\ell}(u)}(W,\mathcal G,\ell)\right\}
\cup
\bigl(\{e_u\}\text{ if }z_u(W,\mathcal G,\ell)=1\bigr),
&g_{\mathcal G,\ell}(u)\in[k].
\end{cases}
\]
For each \(v\in V\), define only the coordinate projection
\[
S_{vu}(W,\mathcal G,\ell)
=\{x_v:x\in S_{\bullet u}(W,\mathcal G,\ell)\}.
\]

%%% FIELD def-completion-handle
For an exactly encoded rational input \(\widetilde W\), its recovered group system \(\mathcal G\), the known labels \(\ell\), an integer \(t\) with \(0\le t\le pq\), and either no prescribed pair or one prescribed pair \((g,r,c)\), define the promise decision problem \(\mathrm{COMP\text{-}EDGE}\) to ask whether there is a legal \(\omega\in\Omega(\widetilde W,\mathcal G,\ell)\), satisfying \(\omega_g=(r,c)\) when prescribed, such that \(E(\omega)\le t\).  The promise is that \((\widetilde W,\mathcal G,\ell)\) is an exactly encoded instance of the canonical Yang model.  The associated optimization and search task is \(\operatorname{CompOpt}(\widetilde W,\mathcal G,\ell;g,r,c)\).  The unrestricted complexity classification remains the open question.  Two auditable handles are retained: exact fraction-free verification of a proposed completion, and decomposition over row--source blocks when \(\widetilde W\) and all candidate groups respect a common block partition.

%%% FIELD proof-observation-group-partition
By \(\text{lem:yang-algorithm-equivalence}\), Algorithm 1 returns the ancestral ordered grouping as a partition in causal group order.  Each substantive label is assigned once and removed once.  The known observation-label map in \(\text{ass:known-observation-labels}\) may remove an inadmissible center role, but it neither duplicates a label nor moves it to a second recovered group.  Hence the legal candidate-center sets \(R_1,\ldots,R_k\) are pairwise disjoint.

By the definition of \(V_{\mathrm{fix}}\),
\[
V_{\mathrm{fix}}(\mathcal G,\ell)
=V\setminus\bigcup_{g=1}^kR_g.
\]
Therefore
\[
V=\left(\bigsqcup_{g=1}^kR_g\right)
\sqcup V_{\mathrm{fix}}(\mathcal G,\ell).
\]
For every \(u\in V\), exactly one branch in the definition of \(g_{\mathcal G,\ell}\) consequently applies: either \(u\in R_g\) for a unique \(g\in[k]\), or \(u\in V_{\mathrm{fix}}\) and the value is \(\bot\).

A legal completion \(\omega=((r_h,c_h))_{h=1}^k\) selects exactly one center \(r_g\in R_g\).  Thus, when \(u\in R_g\), the variable \(u\) is the cogent center exactly when \(r_g=u\); otherwise it has the mleaf role in that reconstruction.  By the construction of \(V_{\mathrm{fix}}\) and the definition of \(\Omega\), every label in \(V_{\mathrm{fix}}\) remains a fixed singleton mleaf.  This proves every assertion.

%%% FIELD proof-normalized-reconstruction
Fix \(\omega\in\Omega(W,\mathcal G,\ell)\), and abbreviate \(R=R(\omega)\), \(C=C(\omega)\), \(L=L(\omega)\), and \(H=H(\omega)\).  The pivot and causal-order conditions included in legality of \(\omega\) give
\[
W_{r_gc_g}\ne0\quad(g\in[k]),
\qquad
W_{r_gc_h}=0\quad(g<h).
\]
Thus \(\Delta_\omega\) is invertible, and
\[
(Q_\omega)_{gg}=1,
\qquad
(Q_\omega)_{gh}=0\quad(g<h).
\]
Therefore \(Q_\omega\) is unit lower triangular and invertible.

Put \(\overline N_C=\Delta_\omega N_C\).  In row order \((R,L)\), the reconstructed equations are
\[
X_R=A_\omega X_R+B^C_\omega N_H+\overline N_C,
\qquad
X_L=D_\omega X_R+B^L_\omega N_H.
\]
Since \(I-A_\omega=Q_\omega^{-1}\), the mixing matrix for source order \((\overline N_C,N_H)\) is
\[
\overline W_\omega=
\begin{pmatrix}
Q_\omega&Q_\omega B^C_\omega\\
D_\omega Q_\omega&D_\omega Q_\omega B^C_\omega+B^L_\omega
\end{pmatrix}.
\]
The four defining identities in \(\text{def:normalized-reconstruction}\) give, entry by entry,
\[
Q_\omega=W[R,C]\Delta_\omega^{-1},
\qquad
Q_\omega B^C_\omega=W[R,H],
\]
\[
D_\omega Q_\omega=W[L,C]\Delta_\omega^{-1},
\qquad
D_\omega Q_\omega B^C_\omega+B^L_\omega=W[L,H].
\]
Consequently
\[
\overline W_\omega=
\begin{pmatrix}
W[R,C]\Delta_\omega^{-1}&W[R,H]\\
W[L,C]\Delta_\omega^{-1}&W[L,H]
\end{pmatrix}.
\]
Returning from \(\overline N_C\) to \(N_C\) multiplies the selected columns by \(\Delta_\omega\), and hence recovers exactly \(W\).

The inverse of a unit lower-triangular matrix is unit lower triangular, so \(A_\omega=I-Q_\omega^{-1}\) is strictly lower triangular.  The center submodel is therefore acyclic.  The remaining blocks have the canonical types: \(D_\omega\) points from centers to mleafs, \(B^C_\omega\) and \(B^L_\omega\) load root latent sources, and mleafs have no substantive children.  All pivot, label, latent-child, and canonical support conditions are among the constraints defining \(\Omega\); hence \(M_\omega\) is a legal canonical reconstruction.

Finally let \(S\) be an invertible diagonal column-scaling matrix and put \(W'=WS\).  With the induced selected and complementary diagonal blocks,
\[
\Delta'_\omega=\Delta_\omega S_C,
\qquad
W'[R,C](\Delta'_\omega)^{-1}=W[R,C]\Delta_\omega^{-1}=Q_\omega.
\]
Thus \(A'_\omega=A_\omega\) and \(D'_\omega=D_\omega\), while
\[
(B^C_\omega)'=B^C_\omega S_H,
\qquad
(B^L_\omega)'=B^L_\omega S_H.
\]
The latter changes only rescale the independent latent sources.  Column permutations merely relabel \(C\) and \(H\).  Hence the reconstruction is invariant under the recovered mixing indeterminacy.

%%% FIELD proof-local-effect-column
Fix a legal \(\omega\).  In substantive row order \((R(\omega),L(\omega))\), the coefficient matrix among substantive variables is
\[
G_\omega=
\begin{pmatrix}
A_\omega&0\\
D_\omega&0
\end{pmatrix}.
\]
The zero right-hand blocks follow from the canonical requirement that an mleaf has no substantive children.  Since \(I-A_\omega=Q_\omega^{-1}\), direct block multiplication yields
\[
(I-G_\omega)^{-1}
=
\begin{pmatrix}
Q_\omega&0\\
D_\omega Q_\omega&I
\end{pmatrix}.
\]

For the acyclic linear structural equations \(X=G_\omega X+\eta\), an additive intervention of size \(t\) in the equation for \(u\) changes the solution by \(t(I-G_\omega)^{-1}e_u\).  Differentiating at any \(t\) shows that the causal total-effect column is the \(u\)-th column of \((I-G_\omega)^{-1}\).  This derives the causal target from the structural equations rather than defining it from \(W\).

If \(u\in L(\omega)\), that column is \(e_u\).  Hence the diagonal effect is one and every off-diagonal effect is zero.  By \(\text{lem:observation-group-partition}\), this covers both \(g_{\mathcal G,\ell}(u)=\bot\) and the case \(u\in R_g\) with \(r_g\ne u\).

If \(u=r_g\), then the same column is column \(g\) of
\[
\begin{pmatrix}Q_\omega\\D_\omega Q_\omega\end{pmatrix}
=W[V,C(\omega)]\Delta_\omega^{-1}.
\]
It is therefore
\[
\frac{W_{\bullet c_g}}{W_{uc_g}}.
\]
The denominator is nonzero by the legal-pivot condition.  Taking coordinate \(v\) proves every displayed scalar case in the theorem.  The vector depends only on \((r_g,c_g)\), so choices outside the unique group containing \(u\) do not alter the entire intervention column.

%%% FIELD stmt-deterministic-joint-image
For every fixed real \(W\), recovered group system \(\mathcal G(W,\ell)\), known label map \(\ell\), and nonempty finite set \(\Omega^{\star}(W,\mathcal G,\ell)\) of globally minimum legal completions, the exact joint intervention-column image is
\[
\{T_{\omega,\bullet u}:\omega\in\Omega^{\star}(W,\mathcal G,\ell)\}
=S_{\bullet u}(W,\mathcal G,\ell)
\qquad(u\in V).
\]
If \(g_{\mathcal G,\ell}(u)=\bot\), this set is \(\{e_u\}\).  If \(g_{\mathcal G,\ell}(u)=g\in[k]\), it is
\[
\left\{\frac{W_{\bullet c}}{W_{uc}}:(u,c)\in P_g(W,\mathcal G,\ell)\right\}
\cup
\bigl(\{e_u\}\text{ if some }(r,c)\in P_g(W,\mathcal G,\ell)\text{ has }r\ne u\bigr).
\]
Every displayed vector is attained by one \(\omega\in\Omega^{\star}\), and
\[
|S_{\bullet u}|\le |K_g|+1\le q+1
\]
on the cogent-group branch.  Moreover \(S_{vu}=\{x_v:x\in S_{\bullet u}\}\) for every \(v\in V\).  No factorization of the full total-effect-matrix image across groups is implied by the marginal projections \(P_g\).

%%% FIELD proof-deterministic-joint-image
Fix \(u\in V\).  If \(g_{\mathcal G,\ell}(u)=\bot\), \(\text{lem:observation-group-partition}\) makes \(u\) a fixed mleaf, and \(\text{thm:local-effect-column}\) gives \(T_{\omega,\bullet u}=e_u\) for every \(\omega\in\Omega^\star\).  This is the first branch of \(\text{def:effect-image}\).

Now suppose \(g_{\mathcal G,\ell}(u)=g\in[k]\).  Partition the finite argmin into
\[
\Omega_u^\star=\{\omega\in\Omega^\star:r_g=u\},
\qquad
\Omega_{\neg u}^\star=\{\omega\in\Omega^\star:r_g\ne u\}.
\]
For \(\omega\in\Omega_u^\star\), the definition of \(P_g\) gives \((u,c_g)\in P_g\), and the local-effect theorem gives
\[
T_{\omega,\bullet u}=\frac{W_{\bullet c_g}}{W_{uc_g}}.
\]
Conversely, if \((u,c)\in P_g\), the defining existential quantifier for \(P_g\) supplies a single \(\omega_c\in\Omega^\star\) with \(\omega_{c,g}=(u,c)\); the same local identity makes \(\omega_c\) attain the whole normalized vector, not merely its coordinates separately.  Each denominator is nonzero because \((u,c)\) is a legal pivot.

For \(\omega\in\Omega_{\neg u}^\star\), the partition lemma makes \(u\) an mleaf and the local theorem gives \(T_{\omega,\bullet u}=e_u\).  Such an \(\omega\) exists exactly when some projected pair \((r,c)\in P_g\) has \(r\ne u\), which is exactly the event \(z_u=1\).  The two parts of the partition therefore prove the asserted set equality and attainment.

There are at most \(|K_g|\) normalized columns and at most one additional vector \(e_u\).  Since \(K_g\subseteq C_{\mathrm{src}}\) and \(|C_{\mathrm{src}}|=q\),
\[
|S_{\bullet u}|\le |K_g|+1\le q+1.
\]
Taking the \(v\)-th coordinate of each vector is exactly the definition of \(S_{vu}\), proving the coordinate claim.

Finally, each \(P_g\) records only whether a local pair occurs in some global minimizer.  For pairs from two distinct groups, the two existential witnesses may be different and need not coexist in one \(\omega\).  Hence the proof yields no Cartesian factorization of \(\{T_\omega:\omega\in\Omega^\star\}\); the exact full-matrix image remains indexed by the joint set \(\Omega^\star\).

%%% FIELD stmt-observational-completeness
For every \(\theta\in\Theta_{p,q,\ell}^{\mathrm{gen}}\) and every \((\theta,M,P_{\mathrm{obs}},W,\mathcal G(W,\ell))\in\mathfrak M_{p,q,\ell}^{\mathrm{Yang}}(\mathbb R)\), the minimum reconstructions \(\{M_\omega:\omega\in\Omega^\star(W,\mathcal G(W,\ell),\ell)\}\) are exactly the observationally compatible minimum-edge Yang DOG class, up to source permutation and nonzero scaling.  All members satisfy Assumption 3 simultaneously on this probability-one domain.  Consequently the deterministic set \(S_{\bullet u}(W,\mathcal G(W,\ell),\ell)\) is the complete observational joint image of the labeled total-effect column.  Each vector is attained, and two unequal vectors are attained by faithful observationally equivalent countermodels.

%%% FIELD proof-observational-completeness
By \(\text{ass:separability}\), \(P_{\mathrm{obs}}\) identifies the retained mixing matrix up to column permutation and nonzero scaling.  By \(\text{lem:normalized-reconstruction}\), the normalization, reconstructed substantive coefficients, and normalized columns are invariant under precisely this indeterminacy.  By the revalidated \(\text{lem:yang-algorithm-equivalence}\), under \(\text{ass:canonical-lvsemme}\), \(\text{ass:minimality}\), \(\text{ass:rank-faithfulness}\), and \(\text{ass:known-observation-labels}\), Algorithm 2 enumerates the legal AOG reconstructions and its minimum-edge members are the Yang DOG class on the domain of Theorem 3.  The original model supplies a legal choice; thus \(\Omega\ne\varnothing\).  Since \(\Omega\) is finite, \(\Omega^\star\ne\varnothing\).

The additional probability-one membership claim is not inferred from the deterministic algebra.  It is exactly the finite-class conclusion of \(\text{lem:yang-dog-faithfulness-preservation}\), which applies on \(\Theta_{p,q,\ell}^{\mathrm{gen}}\): the parameter rewrite to every member is invertible, each exceptional faithfulness locus is null, and the finite union of those loci remains null.  Thus every attaining reconstruction used below remains in the published rank-faithful class simultaneously.

Apply \(\text{thm:deterministic-joint-effect-image}\) to the now identified finite class.  It gives the complete joint column image and an attaining minimum reconstruction for every vector.  All reconstructed retained mixing matrices equal \(W\) after the permitted source rescalings.  Restore the removed one-hot measurement-error columns identically, and rescale or permute the corresponding independent source laws inversely.  The resulting models induce the same \(P_{\mathrm{obs}}\).  Hence, if \(x\ne y\) are in the joint image, their attaining reconstructions are faithful, minimum-edge, observationally equivalent countermodels with unequal labeled total-effect columns.

The qualifier is relative Zariski generic within each fixed real canonical coefficient stratum.  By the definition of \(\Theta_{p,q,\ell}^{\mathrm{gen}}\), the deleted coefficient sets are proper algebraic subsets and hence Lebesgue-null in that stratum's free Euclidean coordinates.  Separability is a distinct source-law assumption and is not included in this Zariski assertion.  An exactly encoded rational input is deterministic, so no probability statement about rational points is made.

%%% FIELD proof-complete-effect-image
Fix a tuple covered by the theorem.  By \(\text{thm:observational-dog-column-completeness}\), the observationally compatible minimum-edge DOG class is indexed by \(\Omega^\star\), and by \(\text{thm:deterministic-joint-effect-image}\) its exact joint \(u\)-intervention-column image is \(S_{\bullet u}\).  The revised \(\text{def:effect-image}\) defines
\[
S_{vu}=\{x_v:x\in S_{\bullet u}\}.
\]
Therefore
\[
S_{vu}
=\{\tau_\omega(v\leftarrow u):\omega\in\Omega^\star\}.
\]
This equality derives the observable functional from the structural total effects; it does not define the causal target to equal an observed quantity.

If \(g_{\mathcal G,\ell}(u)=\bot\), the joint set is \(\{e_u\}\), giving \(S_{uu}=\{1\}\) and \(S_{vu}=\{0\}\) for \(v\ne u\).  On the cogent branch, taking coordinate \(v\) in the joint formula gives
\[
S_{vu}
=\left\{\frac{W_{vc}}{W_{uc}}:c\in L_u\right\}
\cup\bigl(\{0\}\text{ if }z_u=1\bigr)
\quad(v\ne u),
\]
and gives \(S_{uu}=\{1\}\).  The denominators are nonzero legal pivots.  Projection cannot increase cardinality, so
\[
|S_{vu}|\le |S_{\bullet u}|\le q+1.
\]

Every scalar value is the \(v\)-coordinate of at least one vector in the finite joint image.  The joint-image theorem supplies one minimum completion attaining that vector, hence attaining the scalar.  If two scalar values differ, their attaining vectors differ, and \(\text{thm:observational-dog-column-completeness}\) supplies faithful observationally equivalent countermodels.  Its final paragraph also establishes exactly the stated relative-Zariski, Lebesgue-almost-everywhere, source-law, and rational-input qualifications.

%%% FIELD stmt-joint-singleton
For every fixed instance covered by \(\text{thm:deterministic-joint-effect-image}\), the joint set \(S_{\bullet u}\) is a singleton if and only if either \(g_{\mathcal G,\ell}(u)=\bot\), or \(g_{\mathcal G,\ell}(u)=g\in[k]\) and one of the following holds: \(L_u=\varnothing\); \(z_u=1\) and \(W_{vc}=0\) for every \(c\in L_u\) and every \(v\ne u\); or \(z_u=0\), \(L_u\ne\varnothing\), and
\[
W_{vc}W_{ud}-W_{vd}W_{uc}=0
\]
for every \(v\in V\) and every \(c,d\in L_u\).  The common vector is \(e_u\) in the first two cogent cases and any normalized column \(W_{\bullet c}/W_{uc}\) in the last case.

%%% FIELD proof-joint-singleton
The fixed-mleaf branch is \(\{e_u\}\) by \(\text{def:effect-image}\).  Now let \(g_{\mathcal G,\ell}(u)=g\in[k]\).  Every \(c\in L_u\) is paired with \(u\) in a legal completion, so
\[
W_{uc}\ne0\qquad(c\in L_u).
\]
If \(L_u=\varnothing\), nonemptiness of \(P_g\) forces every attaining projected pair to have center different from \(u\), so \(z_u=1\) and \(S_{\bullet u}=\{e_u\}\).

Suppose \(z_u=1\).  Then \(e_u\in S_{\bullet u}\).  A normalized column equals \(e_u\) exactly when, for every \(v\ne u\),
\[
\frac{W_{vc}}{W_{uc}}=0
\quad\Longleftrightarrow\quad
W_{vc}=0,
\]
where division is justified by the legal-pivot condition.  Thus the joint set is a singleton exactly under the stated zero tests, and its value is \(e_u\).

Finally suppose \(z_u=0\) and \(L_u\ne\varnothing\).  There is no adjoined \(e_u\), so singletonness is equality of all normalized column vectors.  For \(c,d\in L_u\), their \(v\)-coordinates are equal exactly when
\[
\frac{W_{vc}}{W_{uc}}=\frac{W_{vd}}{W_{ud}}
\quad\Longleftrightarrow\quad
W_{vc}W_{ud}-W_{vd}W_{uc}=0,
\]
again because both denominators are nonzero.  Requiring this for every \(v\in V\) is equivalent to vector equality.  The three cases exhaust the possibilities and prove both directions.

%%% FIELD proof-singleton-minor
By \(\text{thm:complete-effect-image}\), \(S_{vu}\) is the \(v\)-coordinate projection of the exact joint image.  If \(v=u\), every normalized column and \(e_u\) have coordinate one, so \(S_{uu}=\{1\}\).  If \(v\ne u\) and \(g_{\mathcal G,\ell}(u)=\bot\), then \(S_{vu}=\{0\}\).

Let \(v\ne u\) and \(g_{\mathcal G,\ell}(u)=g\in[k]\).  For every \(c\in L_u\), legality gives \(W_{uc}\ne0\).  If \(L_u=\varnothing\), nonemptiness of the finite argmin gives an attaining pair with center different from \(u\), so \(z_u=1\) and \(S_{vu}=\{0\}\).

If \(z_u=1\), zero is in the scalar projection.  It is a singleton exactly when each displayed ratio is zero, and
\[
\frac{W_{vc}}{W_{uc}}=0
\quad\Longleftrightarrow\quad
W_{vc}=0.
\]
This proves the zero branch, including empty \(L_u\) vacuously.

If \(z_u=0\) and \(L_u\ne\varnothing\), singletonness means that all scalar ratios agree.  For \(c,d\in L_u\),
\[
\frac{W_{vc}}{W_{uc}}=\frac{W_{vd}}{W_{ud}}
\quad\Longleftrightarrow\quad
W_{vc}W_{ud}-W_{vd}W_{uc}=0.
\]
Thus the pairwise minors are necessary and sufficient, and any admissible ratio is the common value.  These cases are exhaustive.

%%% FIELD stmt-sharpness
For every integer \(q\ge1\), the bound \(q+1\) in \(\text{thm:deterministic-joint-effect-image}\) is attained by a separable canonical, minimal, LV-SEM-ME rank-faithful Yang model with known labels: there are two measured substantive labels \(u\ne v\), one recovered group \(R_1=\{u,v\}\), \(K_1=C_{\mathrm{src}}=[q]\), and an exactly encoded rational retained mixing representative for which
\[
|S_{\bullet u}(W,\mathcal G(W,\ell),\ell)|=q+1.
\]
Every one of the \(2q\) legal row--column choices is a global minimum and has \(2q-1\) structural edges.

%%% FIELD proof-sharpness
Fix \(q\ge1\).  Consider a two-row matrix with columns
\[
w_j=\begin{pmatrix}a_j\\b_j\end{pmatrix}
\qquad(j\in[q]),
\]
where every \(a_j\) and \(b_j\) is nonzero and the ratios \(b_j/a_j\) are pairwise distinct.  These requirements are the nonvanishing of finitely many coordinate polynomials and minors \(a_jb_h-a_hb_j\).

We also require the coefficients to avoid Yang et al.'s finite collection of proper algebraic faithfulness-exceptional loci for this fixed canonical stratum and its finite DOG class.  Such a rational choice exists.  Indeed, the product \(F\) of the nonzero polynomials defining the coordinate, minor, and exceptional conditions is a nonzero real polynomial.  We prove directly that \(\mathbb Q^d\) is Zariski dense in \(\mathbb R^d\): for \(d=1\), a nonzero polynomial has only finitely many roots; inductively, write a nonzero polynomial as a polynomial in its last variable, choose rational values of the first \(d-1\) variables at which a nonzero coefficient remains nonzero, and then choose a rational last coordinate outside the finite root set.  Applying this argument to \(F\) gives rational \((a_1,b_1,\ldots,a_q,b_q)\) with \(F\ne0\).

Let both substantive labels be measured, let \(R_1=\{u,v\}\), and let \(K_1=[q]\).  Select a center \(r\in\{u,v\}\) and a source column \(c\in[q]\), and denote the other row by \(l\).  There is one center, so
\[
Q_\omega=(1),
\qquad
A_\omega=(0),
\qquad
D_\omega=\frac{W_{lc}}{W_{rc}}\ne0.
\]
For every complementary column \(h\ne c\),
\[
(B^C_\omega)_h=W_{rh}\ne0,
\qquad
(B^L_\omega)_h
=W_{lh}-\frac{W_{lc}}{W_{rc}}W_{rh}
=\frac{W_{lh}W_{rc}-W_{lc}W_{rh}}{W_{rc}}\ne0.
\]
The last inequality is exactly the pairwise-ratio condition, and the denominator is a nonzero pivot.  Thus every row--column choice is canonical: there is one nonzero center-to-mleaf edge, and each of the \(q-1\) complementary root latent sources has nonzero loadings to both substantive variables.  Its edge count is
\[
E(\omega)=0+(q-1)+1+(q-1)=2q-1.
\]
There are no other groups or choices, so all \(2q\) choices are legal global minimizers.

Minimality follows from \(\text{lem:yang-minimality-criterion}\).  In any reconstruction, each root latent source has exactly the center \(r\) and the mleaf \(l\) as substantive children, and the nonzero edge \(r\to l\) gives
\[
r\in\operatorname{An}(l),
\qquad
r\notin\operatorname{An}(r),
\]
where \(\operatorname{An}\) denotes strict ancestors.  Hence \(\operatorname{An}(l)\nsubseteq\operatorname{An}(r)\), so the cited nonminimality witness does not exist.  The construction is minimal.  By the choice outside the exceptional loci and \(\text{lem:yang-dog-faithfulness-preservation}\), Assumption 3 holds simultaneously for all \(2q\) reconstructions.  Attach independent source laws in the separability domain of \(\text{ass:separability}\), and append the two one-hot measurement-error sources dictated by the known labels.  This places the example in the stated published class.

When the chosen center is \(u\), the \(q\) column choices give
\[
\frac{W_{\bullet j}}{W_{uj}}
=\begin{pmatrix}1\\b_j/a_j\end{pmatrix},
\qquad j\in[q],
\]
which are pairwise distinct.  When the center is \(v\), the label \(u\) is an mleaf and its effect column is \(e_u=(1,0)^\mathsf T\).  Since every \(b_j/a_j\ne0\), this vector is distinct from the other \(q\) vectors.  Therefore \(|S_{\bullet u}|=q+1\), proving sharpness.

%%% FIELD proof-projected-oracle
Assume the polynomial-time solver \(\mathsf A_{\min}\).  One unconstrained call returns
\[
e^\star=\min_{\omega\in\Omega}E(\omega)
\]
and an attaining completion.  For each \(g\in[k]\) and \((r,c)\in R_g\times K_g\), make a constrained call and write its optimum as \(e_{g,r,c}\), with value \(+\infty\) if the call certifies infeasibility.  The constrained feasible set is contained in \(\Omega\), so \(e_{g,r,c}\ge e^\star\), and the definition of \(P_g\) gives the exact equivalence
\[
(r,c)\in P_g
\quad\Longleftrightarrow\quad
e_{g,r,c}=e^\star.
\]
The forward implication uses the global minimizer witnessing membership in \(P_g\); the reverse implication uses the completion returned by the constrained call.  Retain that witness whenever equality holds.

Because the row sets \(R_g\) are disjoint, the products \(R_g\times K_g\) are disjoint subsets of \(V\times C_{\mathrm{src}}\).  Hence
\[
\sum_{g=1}^k|R_g||K_g|
\le |V|\,|C_{\mathrm{src}}|=pq.
\]
Thus at most \(pq\) constrained calls compute all projections and an attaining global minimizer for every retained pair.

For each \(u\), collect the pairs with first coordinate \(u\) and check whether another first coordinate occurs in its unique group.  The revised \(\text{def:effect-image}\) then forms the exact joint set \(S_{\bullet u}\), and its coordinate projections give every \(S_{vu}\).  At most \(q\) normalized columns are processed for each of \(p\) interventions and \(p\) coordinates, requiring \(O(p^2q)\) additional rational operations.  Equality uses numerator-zero tests and the exact cross-products
\[
W_{vc}W_{ud}=W_{vd}W_{uc},
\]
so no numerical division is required for certification.

If a scalar image is a singleton, these comparisons return its common value.  If it is not, choose two unequal generated coordinates.  A normalized-column value uses the stored minimizer for its pair, and an adjoined zero uses any stored alternative-center minimizer.  The complete-effect theorem proves attainment and supplies the faithful observationally equivalent countermodel interpretation on the generic population domain.  Under the assumed polynomial bit-time oracle, the bounded number of calls and exact rational postprocessing are polynomial.

%%% FIELD proof-no-measurement
On the stated subclass every substantive variable is directly observed and no substantive row can take an mleaf role.  The known-label restriction and \(\text{lem:observation-group-partition}\) therefore give singleton center-row sets \(R_g=\{u_g\}\), no fixed-mleaf labels, and \(z_u=0\) for every \(u\).  The joint image becomes
\[
S_{\bullet u}(W,\mathcal G(W,\ell),\ell)
=\left\{\frac{W_{\bullet c}}{W_{uc}}:c\in L_u(W,\mathcal G(W,\ell),\ell)\right\},
\]
and hence
\[
S_{vu}(W,\mathcal G(W,\ell),\ell)
=\left\{\frac{W_{vc}}{W_{uc}}:c\in L_u(W,\mathcal G(W,\ell),\ell)\right\}.
\]
Each denominator is a legal nonzero pivot.  Normalizing candidate source column \(c\) to have \(u\)-entry one gives exactly this vector and this \(v\)-coordinate.  With no measured-variable center exchange and no mleaf rows, the remaining ambiguity is precisely the assignment of an observed structural-noise column versus a root-latent column in canonical latent-variable LiNGAM.  Thus this is the normalized mixing-column effect set used in that model.

On the shared generic domain, \(\text{thm:complete-effect-image}\) says that the displayed coordinate set is the complete observational effect image.  Therefore its \(v\)-coordinate is a singleton exactly when the total effect of \(u\) on \(v\) is generically identified.  The revalidated \(\text{lem:tramontano-unknown-graph-total-effect-criterion}\) states that, for the same canonical latent-variable LiNGAM observational class and the same intervention derivative, this occurs exactly under its Theorem 3.3 graphical condition.  Chaining the two equivalences proves the proposition.  The scalar minor test is the coordinatewise algebraic form of the same normalized-column equality.

%%% FIELD stmt-comp-edge-np
On the promised exactly encoded rational Yang inputs, \(\mathrm{COMP\text{-}EDGE}\), both unconstrained and with one prescribed local pair, belongs to promise-\(\mathrm{NP}\).  A legal completion \(\omega\) is a polynomial-size certificate, and its legality and the inequality \(E(\omega)\le t\) are verifiable in polynomial bit time by fraction-free exact reconstruction and support counting.  Since \(0\le E(\omega)\le pq\), exact optimization is in \(\mathrm{FP}^{\mathrm{NP}}\), and an attaining completion can also be recovered in \(\mathrm{FP}^{\mathrm{NP}}\) by prefix self-reduction of this verifier.

%%% FIELD proof-comp-edge-np
An assignment \(\omega=((r_g,c_g))_{g=1}^k\) contains at most \(k\le\min(p,q)\) row--column pairs, so its binary encoding has length polynomial in the input length.  Given \(\omega\), first check \(r_g\in R_g\), \(c_g\in K_g\), the optional equality \(\omega_g=(r,c)\), distinctness of the selected rows and columns, the known-label conditions, and every explicitly listed Algorithm-2 pivot and canonicality constraint.  These are polynomially many exact comparisons.  Reject if any pivot \(\widetilde W_{r_gc_g}\) is zero.

Clear the input denominators and compute \(Q_\omega^{-1}\), \(A_\omega\), \(B^C_\omega\), \(D_\omega\), and \(B^L_\omega\) by fraction-free Gaussian elimination and exact matrix multiplication.  Polynomial bit size follows directly from determinant bounds.  For an \(m\times m\) integer matrix whose entries have at most \(B'\) bits, the Leibniz expansion has \(m!\) products, each with at most \(mB'\) bits, so every determinant has
\[
O\bigl(m(B'+\log m)\bigr)
\]
bits.  Cramer's rule expresses every inverse entry as a ratio of two such minors; reducing fractions cannot increase the numerator or denominator beyond polynomial bit length.  The subsequent polynomially many sums and products in the four reconstruction blocks retain polynomial bit length.  Thus fraction-free elimination performs polynomially many operations on polynomial-size integers.

After reconstruction, verify the lower-triangular, acyclicity, latent-child, mleaf, and label support conditions by exact zero tests.  Count the nonzero entries of the four blocks.  Their total number of positions is
\[
k^2+k(q-k)+(p-k)k+(p-k)(q-k)=pq,
\]
so comparison with \(t\) is polynomial.  This verifier accepts exactly the promised yes instances, proving promise-\(\mathrm{NP}\) membership.

Because the integer optimum lies in \([0,pq]\), binary search with an \(\mathrm{NP}\) oracle determines it using \(O(\log(pq+1))\) queries.  To recover a witness, encode the certificate as a polynomial-length bit string and, one bit at a time, ask the same polynomial verifier whether an accepting completion with the already chosen prefix and the candidate next bit exists at the optimal threshold.  At least one branch remains feasible until a full accepting certificate is obtained.  This uses polynomially many \(\mathrm{NP}\) queries and proves both \(\mathrm{FP}^{\mathrm{NP}}\) claims.  It does not classify the unrestricted problem as polynomial-time solvable or hard.

%%% FIELD stmt-block-decomposition
Suppose \(V=\bigsqcup_{j=1}^sV_j\) and \(C_{\mathrm{src}}=\bigsqcup_{j=1}^sC_j\), with \(\widetilde W[V_i,C_j]=0\) whenever \(i\ne j\), and suppose every candidate group \(g\) is contained in one block in the sense that \(R_g\subseteq V_j\) and \(K_g\subseteq C_j\) for a unique \(j\).  Then, after simultaneous block permutations, every legal normalized reconstruction is a direct sum of its block reconstructions, legality factors across blocks, and
\[
\Omega=\prod_{j=1}^s\Omega_j,
\qquad
E(\omega)=\sum_{j=1}^sE_j(\omega_j),
\qquad
e^\star=\sum_{j=1}^se_j^\star,
\qquad
\Omega^\star=\prod_{j=1}^s\Omega_j^\star.
\]
A prescribed pair affects only its containing block.

%%% FIELD proof-block-decomposition
Fix a legal choice \(\omega\).  Because every selected pair \((r_g,c_g)\) lies in one common block, reordering selected rows and columns by block makes \(\Delta_\omega\) block diagonal.  The assumed zero cross-block submatrices make \(\widetilde W[R(\omega),C(\omega)]\) block diagonal as well.  Therefore
\[
Q_\omega
=\widetilde W[R(\omega),C(\omega)]\Delta_\omega^{-1}
=\bigoplus_{j=1}^sQ_{\omega_j},
\qquad
Q_\omega^{-1}=\bigoplus_{j=1}^sQ_{\omega_j}^{-1}.
\]
The same row--column partition makes \(\widetilde W[R,H]\), \(\widetilde W[L,C]\), and \(\widetilde W[L,H]\) block diagonal.  Substitution into \(\text{def:normalized-reconstruction}\) gives
\[
A_\omega=\bigoplus_jA_{\omega_j},
\qquad
B^C_\omega=\bigoplus_jB^C_{\omega_j},
\qquad
D_\omega=\bigoplus_jD_{\omega_j},
\qquad
B^L_\omega=\bigoplus_jB^L_{\omega_j}.
\]

Membership, nonzero-pivot, observation-label, and latent-child conditions are local to their group and hence to its block.  A block direct sum is acyclic exactly when every block is acyclic, and no cross-block coefficient can violate canonicality.  Thus a global assignment is legal exactly when every restriction \(\omega_j\) is legal, proving \(\Omega=\prod_j\Omega_j\).

The support size of a direct sum is the sum of its block support sizes, so \(E(\omega)=\sum_jE_j(\omega_j)\).  Minimizing a sum over a Cartesian product gives
\[
\min_{(\omega_1,\ldots,\omega_s)\in\prod_j\Omega_j}
\sum_jE_j(\omega_j)=\sum_j\min_{\omega_j\in\Omega_j}E_j(\omega_j).
\]
Equality is attained by choosing a minimizer in every block, and a global choice is minimizing only if each component is minimizing.  This proves the formulas for \(e^\star\) and \(\Omega^\star\).  A prescribed pair lies in a unique block and restricts only that factor.

%%% FIELD stmt-block-tractability
Under the hypotheses of \(\text{lem:block-completion-decomposition}\), suppose each block contains at most a fixed number \(b\) of nontrivial groups, where a group is nontrivial when it has more than one legal local row--column pair.  Then both the unconstrained and prescribed-pair versions of \(\operatorname{CompOpt}\) can be solved, with an attaining completion or an infeasibility certificate, in
\[
(pq)^b\operatorname{poly}(p,q,B)
\]
bit operations for input bit length \(B\).  Thus the problem is polynomial-time solvable for fixed \(b\), even when the total number of ambiguous groups over all blocks is unbounded.

%%% FIELD proof-block-tractability
In a fixed block, all groups with a single legal local pair are forced.  By hypothesis there are at most \(b\) remaining groups.  Each has at most \(|R_g||K_g|\le pq\) candidate pairs, so the number of block assignments is at most
\[
\prod_{g\text{ nontrivial in the block}}|R_g||K_g|
\le(pq)^b.
\]
Enumerate these assignments, impose a prescribed pair when it belongs to this block, reject illegal assignments by the verifier in \(\text{thm:comp-edge-promise-np}\), and evaluate the exact support count of every survivor.  The determinant bit bound proved there shows that each evaluation takes \(\operatorname{poly}(p,q,B)\) bit operations.  Retain a least-support assignment, or report infeasibility if none survives.

By \(\text{lem:block-completion-decomposition}\), the global optimum is the sum of the block optima and a global attaining completion is obtained by concatenating the retained block minimizers.  A prescribed pair changes only one block's enumeration.  There are at most \(p+q\) nonempty blocks, which is absorbed by the polynomial factor, so the total running time is
\[
(pq)^b\operatorname{poly}(p,q,B).
\]
For fixed \(b\) this is polynomial.  The number of nontrivial groups may grow with the number of blocks, so the result permits arbitrarily many ambiguous groups overall; only their interaction width within one block is bounded.
